% mdframed's boxes have their lines, their space inside and around, and
% their colors as in the PDF: the default box; a box with custom margins,
% colors, and line width (the box after it keeps the default line); a box
% with only side lines; and a theorem framed by \newmdtheoremenv, whose
% frame ends at the theorem, with a \ref to it.
\documentclass{article}
\usepackage{amsmath}
\usepackage{xcolor}
\usepackage{mdframed}
\newmdtheoremenv{lemma}{Lemma}
\begin{document}
Text before the box.
\begin{mdframed}
A framed paragraph with $x^2 + y^2 = z^2$ in it.

A second paragraph.
\end{mdframed}
Text between.
\begin{mdframed}[innerleftmargin=20pt, innertopmargin=12pt, linecolor=blue,
  backgroundcolor=yellow!20, linewidth=2pt]
A box with its own margins and colors.
\end{mdframed}
\begin{lemma}\label{lem:one}
A framed lemma: $\sum_{i=1}^n i = n(n+1)/2$.
\end{lemma}
Lemma~\ref{lem:one} is framed.
\begin{mdframed}[topline=false, bottomline=false, leftmargin=20pt, rightmargin=20pt]
A box with only its side lines, set in from the text's margins.
\end{mdframed}
\end{document}
